\documentclass[11pt,a4paper]{article}
\usepackage[margin=23mm]{geometry}
\usepackage[T1]{fontenc}
\usepackage{lmodern,microtype,amsmath,amssymb,amsthm,mathtools,enumitem,xcolor,booktabs}
\usepackage[colorlinks=true,linkcolor=blue!45!black,citecolor=blue!45!black,urlcolor=blue!45!black]{hyperref}
\newtheorem{theorem}{Theorem}[section]
\newtheorem{lemma}[theorem]{Lemma}
\newtheorem{proposition}[theorem]{Proposition}
\newtheorem{corollary}[theorem]{Corollary}
\theoremstyle{remark}\newtheorem{remark}[theorem]{Remark}
\setlength{\parskip}{3pt}
\setlist{nosep}
\allowdisplaybreaks[1]
\title{Thresholds, coherent structures, and dynamics\\of critical focusing dispersive systems\\[5pt]
\large Persistent kinetic amplification, conserved occupations,\\and momentum-adapted motion}
\author{Shuxing Yang}
\date{8 October 2026}
\begin{document}
\maketitle
\begin{abstract}
We connect sharp critical concentration geometry to nonlinear evolution
through three results. For real-even positive elliptic homogeneous
dispersion of order $m>2$, global zero-momentum negative-energy
solutions near critical mass satisfy a lower physical-time mean kinetic
bound $\liminf_{T\to\infty}T^{-1}\int_0^T A(t)\,dt
\ge c|E|d^{-1-m/[2(m-1)]}$. The result covers critical local powers
and a class of strictly positive-definite homogeneous Hartree
interactions, without radiality, a spatial moment, a global kinetic
bound, or symbol convexity. Fourier factorization preserves an exterior
negative kinetic term, while intrinsic-time center construction and
weak energy splitting control the small mass tail.
For local multicomponent systems, Gaussian projection of the scalar
sharp inequality gives an explicit coefficient depending on the entire
conserved occupation spectrum. This yields all-time energy bounds for
specified above-critical masses, including $M<\sqrt2M_c$ for two equal
orthogonal occupations in the two-dimensional quartic system, and an
exact scalar reduction for every unbounded critical-mass evolution.
Finally, near a quartic ground wave with nondegenerate real linearization,
momentum-adapted modulation improves the mean bound to
$c_Q|E|d^{-7/3}$ at zero momentum and to $c_Q|E|d^{-2}$ at arbitrary
momentum. Complete supporting proofs and the remaining dynamical
obstruction are included.
\end{abstract}

The stronger modulation results
retain the explicit assumption $\ker L_+=\operatorname{span}\{\partial_jQ\}$.
The general mean-bound theorem states its conservative strong solution
class; a concrete local-flow proof is supplied for local powers of
order $m=2k\ge4$ in dimensions one through three. None of the results
proves that all nonradial negative-energy solutions blow up in finite time.

\section*{The mathematical connection}
The main question is how critical threshold geometry restricts a
time-dependent solution. Write $M=\|u\|_2^2$ and
$A=\|a(D)^{1/2}u\|_2^2$. Sharp interpolation at critical mass $M_c$
gives only $A(t)\ge2|E|/d$ when $E<0$. We show that nonlinear evolution
forces a substantially higher \emph{long-time average}. The mass surplus
leaves at most $O(d)$ mass outside a persistent core. Quantitative phase
control bounds the motion of that core, and the favorable exterior
kinetic contribution absorbs the tail current. This produces
\[
 \liminf_{T\to\infty}\frac1T\int_0^T A(t)\,dt
 \ge c|E|d^{-\alpha_m},\qquad
 \alpha_m=1+\frac{m}{2(m-1)}.
\]
For amplitude perturbations $u_d(0)=(1+d)^{1/[2(\kappa-1)]}Q$,
where $M(Q)=M_c$ and $\kappa=1+m/N$ locally or $\kappa=2$ for
Hartree, this forces mean amplification at least
$cd^{-m/[2(m-1)]}$. Every prescribed finite factor is crossed at
arbitrarily late times by a global solution for sufficiently small $d$.

The multicomponent result uses the same threshold geometry in a
different dynamical direction. The internal correlation operator
$\Gamma=\int|U(x)\rangle\langle U(x)|\,dx$ is conserved.
An explicit occupation-sensitive sharp-scalar inequality turns that
conservation law into a uniform bound for all times. It also makes
the assertion of scalar behavior at threshold exact for the entire
evolution, not just for a concentrating subsequence.

For a nondegenerate quartic ground component, the inverse phase
operator defines dual translation coordinates in which conserved
momentum determines the leading drift. At zero momentum the drift
becomes quadratic in the perturbation, improving the exponent. At
nonzero momentum the signed main term in moving virial is favorable,
so a weaker improvement survives without imposing $P=0$.

\paragraph{Reading guide.}
Section~\ref{sec:r13-occupations} gives explicit conserved-spectrum
bounds. Section~\ref{sec:r13-dispersion-structure} establishes the
Fourier and interaction estimates and the concrete local flow.
Section~\ref{sec:r13-persistent-general} proves the general mean
bound, and Section~\ref{sec:r13-center} proves the conditional
momentum refinement. The appendices supply scalar compactness,
phase coercivity, quantitative polarization, and an exact obstruction
to a stronger instantaneous current bound.

\section{Conserved occupations and explicit critical bounds}
\label{sec:r13-occupations}

The local critical nonlinearity admits an explicit connection between
scalar sharp interpolation and the entire conserved occupation spectrum.
The auxiliary Gaussian variables below occur only in the proof; the
initial data and the evolution remain deterministic.

Let $\mathcal H$ be a separable complex Hilbert space and suppose the
scalar inequality is
\begin{equation}\label{eq:r13-scalarGN}
 \int |f|^{2(1+r)}\le (1+r)M_c^{-r}M(f)^r A(f),
 \qquad r=m/N>0.
\end{equation}
The positive scalar Fourier multiplier defining $A$ may be anisotropic.
For nonzero $U\in H^{m/2}(\mathbb R^N;\mathcal H)$ set
\[
 \mathcal F(U)=\int\|U\|_{\mathcal H}^{2(1+r)},\qquad
 \Gamma=\int|U(x)\rangle\langle U(x)|\,dx,\qquad
 \Gamma_A=\int a(\xi)|\widehat U(\xi)\rangle
                              \langle\widehat U(\xi)|\,d\xi.
\]
Write $M=\operatorname{Tr}\Gamma$, $A=\operatorname{Tr}\Gamma_A$,
$\mu=\|\Gamma\|_{\rm op}$, and list the positive eigenvalues of
$\Gamma$ as $\lambda_1=\mu\ge\lambda_2\ge\cdots$.
Euler's gamma function is denoted by $\mathrm G$ in this section, to
distinguish it from the correlation operator.

\begin{theorem}[Explicit bound from the occupation spectrum]
\label{thm:r13-Gaussian}
Let $\xi_j$ be independent exponential random variables of mean one.
Define
\begin{equation}\label{eq:r13-Bdef}
 \mathcal B_r(\Gamma)=\frac{1}{\mathrm G(r+2)}
   \mathbb E\left[\xi_1\left(\sum_j\lambda_j\xi_j\right)^r\right].
\end{equation}
Then
\begin{equation}\label{eq:r13-spectralGN}
 \mathcal F(U)\le(1+r)M_c^{-r}\mathcal B_r(\Gamma)A(U).
\end{equation}
The coefficient depends only on the spectrum, is homogeneous of degree
$r$, and satisfies $\mathcal B_r(\Gamma)\le M^r$. Equality in this last
inequality holds precisely when $\Gamma$ has rank one.
\end{theorem}
\begin{proof}
First suppose the internal rank is finite and diagonalize $\Gamma$.
For independent standard circular complex Gaussian variables $z_j$, put
$f_z=\sum_jz_ju_j$. Pointwise, $f_z$ has Gaussian variance
$\rho_U=\|U\|^2$, so
\[
 \mathbb E\mathcal F(f_z)=\mathrm G(r+2)\mathcal F(U).
\]
In the chosen basis $M(f_z)=S:=\sum_j\lambda_j|z_j|^2$.
Apply \eqref{eq:r13-scalarGN} and average. Averaging over the phases
of the $z_j$ removes off-diagonal entries of the kinetic Gram matrix,
giving
\begin{equation}\label{eq:r13-weightedGaussian}
 \mathrm G(r+2)\mathcal F(U)
 \le(1+r)M_c^{-r}\sum_j A(u_j)b_j,
 \qquad b_j=\mathbb E[\xi_j S^r].
\end{equation}
If $\lambda_i\ge\lambda_j$, exchange $\xi_i$ and $\xi_j$ in one
copy of this expectation. With $S'$ denoting the exchanged sum,
\[
 b_i-b_j=\tfrac12\mathbb E[(\xi_i-\xi_j)(S^r-(S')^r)]\ge0,
\]
because $S-S'=(\lambda_i-\lambda_j)(\xi_i-\xi_j)$.
Thus $b_1$ is maximal, proving \eqref{eq:r13-spectralGN}.
H\"older's inequality and convex Jensen at exponent $r+1$ give
\[
 \mathbb E[\xi_1S^r]
 \le(\mathbb E\xi_1^{r+1})^{1/(r+1)}
       (\mathbb ES^{r+1})^{r/(r+1)}
 \le\mathrm G(r+2)M^r.
\]
If a second eigenvalue is positive, $S$ is not proportional to $\xi_1$
almost surely, and H\"older is strict. In rank one the moment is
exactly $M^r\mathrm G(r+2)$.

For infinite rank, truncate in the eigenbasis of $\Gamma$.
The truncated fields converge in $H^{m/2}$ and in the interaction
space. Their moments increase to \eqref{eq:r13-Bdef} by monotone
convergence and remain bounded by $M^r$. This proves the full result.
In particular no identity-covariance Gaussian vector in infinite-dimensional
$\mathcal H$ is assumed to exist.
\end{proof}

\subsection{Explicit quantitative polarization}
For $X=\Gamma/M$ and $Y=\Gamma_A/A$, both positive of trace one, the
normalized distance to a fixed internal direction is exactly
\begin{equation}\label{eq:r13-pdistance}
 \mathfrak p(U):=\inf_{\|\zeta\|=1,q}
 \left\{\frac{M(U-\zeta q)}M+\frac{A(U-\zeta q)}A\right\}
 =2-\lambda_{\max}(X+Y).
\end{equation}
Indeed the pointwise scalar projection minimizes the combined quadratic
norm for each $\zeta$.

\begin{theorem}[Mass and kinetic alignment, $r\ge1$]
\label{thm:r13-explicitpolarization}
For complex Hilbert-valued fields,
\begin{align}
 \frac{\mathcal F(U)}{(1+r)(M/M_c)^r A}
 &\le\frac{1+r\operatorname{Tr}(XY)}{r+1},
 \label{eq:r13-alignment}\\
 \mathfrak p(U)&\le\frac{2(r+1)}r\,\mathfrak d(U),
 \qquad
 \mathfrak d(U)=1-
  \frac{\mathcal F(U)}{(1+r)(M/M_c)^r A}.
 \label{eq:r13-explicitpol}
\end{align}
Consequently
\begin{equation}\label{eq:r13-simpleoccupation}
 \mathcal F(U)\le(1+r)(M/M_c)^r
 \left[1-\frac r{r+1}\left(1-\frac\mu M\right)\right]A.
\end{equation}
For real Hilbert-valued fields with real-even dispersion, the respective factors in
\eqref{eq:r13-alignment} and \eqref{eq:r13-explicitpol} are
\[
 \frac{1+2r\operatorname{Tr}(XY)}{2r+1},
 \qquad \frac{2r+1}{r}.
\]
\end{theorem}
\begin{proof}
For $r\ge1$, Jensen gives
$S^r\le M^{r-1}\sum_i\lambda_i\xi_i^r$. Hence
\[
 \frac{b_j}{\mathrm G(r+2)}
 \le M^{r-1}\left[\lambda_j+\frac{M-\lambda_j}{r+1}\right].
\]
Substitution into \eqref{eq:r13-weightedGaussian} proves
\eqref{eq:r13-alignment}. Since
\[
 \lambda_{\max}(X+Y)
 \ge\frac{\operatorname{Tr}((X+Y)^2)}{\operatorname{Tr}(X+Y)}
 \ge2\operatorname{Tr}(XY),
\]
where the second step is $\operatorname{Tr}((X-Y)^2)\ge0$,
\eqref{eq:r13-pdistance} proves \eqref{eq:r13-explicitpol}.
Also $\operatorname{Tr}(XY)\le\|X\|_{\rm op}=\mu/M$, proving
\eqref{eq:r13-simpleoccupation}.
For real Gaussian projections use $\xi_j=g_j^2$ with standard real
$g_j$. Their moments satisfy
\[
 a_q=\mathbb E\xi_j^q
 =\frac{2^q\mathrm G(q+1/2)}{\sqrt\pi},
 \qquad a_{r+1}=(2r+1)a_r.
\]
The same calculation gives the stated real constants. Finite-rank
approximation proves all assertions in infinite dimension.
\end{proof}

A scalar complex field is a real two-component field
$(\operatorname{Re}u,\operatorname{Im}u)$. Therefore, for $r\ge1$,
there is a constant phase $\theta$ and a real scalar $q$ such that
\begin{equation}\label{eq:r13-explicitphase}
 \frac{M(u-e^{i\theta}q)}{M(u)}+
 \frac{A(u-e^{i\theta}q)}{A(u)}
 \le\frac{2r+1}{r}\mathfrak d(u).
\end{equation}
This supplies the phase-smallness input for near-critical quartic
motion in $N\le4$ directly from sharp scalar interpolation.
The real covariance used for this instantaneous estimate is generally
not conserved. The complex internal covariance in the following
dynamical result is conserved.

\subsection{Consequences for the actual evolution}
Consider the same common scalar system
\[
 i\partial_tU=(a+b)(D)U-\|U\|^{2r}U,\qquad b\ge0,
 \qquad E=\frac{A+B}{2}-\frac{\mathcal F}{2(1+r)}.
\]
Assume its conservative energy-space local flow and continuation
property. All pairwise component inner products, and therefore
$\Gamma$, are conserved. To see this, use the common scalar
self-adjoint Hamiltonian $H_U=(a+b)(D)-\|U\|^{2r}$. For smooth
approximations the equations give
\[
 \frac{d}{dt}\int u_j\overline{u_k}
 =-i\int (H_Uu_j)\overline{u_k}
   +i\int u_j\overline{H_Uu_k}=0.
\]
The quadratic-form or strong-approximation interpretation passes this
identity to the energy-space flow. In infinite internal dimension,
all matrix elements are conserved and determine the positive trace-class
operator $\Gamma$; continuity also follows from
$\|\Gamma(U)-\Gamma(V)\|_1
\le(\|U\|_2+\|V\|_2)\|U-V\|_2$.

\begin{theorem}[Global bounded flow from conserved occupations]
\label{thm:r13-spectrumflow}
If $\mathcal B_r(\Gamma(U_0))<M_c^r$, the solution is global in both
time directions, and
\[
 \sup_t A(U(t))\le
 \frac{2E(U_0)}{1-\mathcal B_r(\Gamma(U_0))/M_c^r},
 \qquad \sup_t B(U(t))\le2E(U_0).
\]
In particular any fixed normalized occupation spectrum of rank greater
than one has an explicitly specified global mass interval extending
strictly above $M_c$. At mass $M_c$, every orbit with unbounded energy
norm lies in one fixed internal line for its entire maximal evolution.
\end{theorem}
\begin{proof}
Theorem~\ref{thm:r13-Gaussian} and conservation of $\Gamma$ yield
\[
 E(U(t))\ge\frac12
 \left[1-\frac{\mathcal B_r(\Gamma(U_0))}{M_c^r}\right]A(U(t))
 +\frac12 B(U(t)).
\]
Conservation and continuation prove the bounds and global existence.
Homogeneity and strictness of $\mathcal B_r(\Gamma)<M^r$ give the
above-critical interval. At $M=M_c$, unboundedness forces rank one.
The fixed range of the conserved $\Gamma$ then contains $U(t,x)$
almost everywhere at every time.
\end{proof}

Proposition~\ref{prop:r13g-local-flow} verifies these flow assumptions
for $m=2k\ge4$ and $N=1,2,3$, also for countably many components.
In particular it applies to quartic dispersion in these dimensions.
No randomness of the initial field,
radial symmetry, or finite spatial moment is required.

The spectral moment is directly computable. For integer $r$, with
$h_r$ the complete homogeneous symmetric polynomial,
\[
 \mathcal B_r(\Gamma)
 =\frac{h_r(\lambda_1,\lambda_1,\lambda_2,\ldots)}{r+1}.
\]
In particular,
\begin{equation}\label{eq:r13-B2}
 \mathcal B_2(\Gamma)=
 \frac{M^2+\operatorname{Tr}\Gamma^2+2M\mu+2\mu^2}{6}.
\end{equation}
For $k$ equal occupations $\lambda_j=M/k$,
\[
 \frac{\mathcal B_r(\Gamma)}{M^r}
 =\frac{\mathrm G(k+r+1)}
 {k^r\mathrm G(k+1)\mathrm G(r+2)}.
\]
Thus in the two-dimensional quartic system, two equal orthogonal
occupations give $\mathcal B_2=M^2/2$: every such $H^2$ datum with
$M<\sqrt2M_c$ has globally bounded evolution.
Formula~\eqref{eq:r13-B2} retains information beyond $M$ and $\mu$.
For $r\ge1$, at fixed $M$ and $\lambda_1$, averaging any two of the
remaining occupations cannot increase $\mathcal B_r$, by convexity
and permutation symmetry of its defining expectation.

These estimates retain the sharp scalar endpoint, but no assertion
of optimality at each higher-rank spectrum is made. Evaluating that
optimal coefficient is a sharp orthonormal kinetic problem. The
Gaussian argument above applies to local powers; Hartree interactions
produce an additional exchange term under Gaussian averaging, so no
Hartree extension is asserted here.

\section{Exterior energy for homogeneous dispersion and attraction}
\label{sec:r13-dispersion-structure}

Fix $m>2$, $s=m/2$, and a real even positive elliptic symbol $a$,
$m$-homogeneous and $C^{N+2}$ away from zero:
\[
 c_a|\xi|^m\le a(\xi)\le C_a|\xi|^m.
\]
For $u\in H^s(\mathbb R^N)$ write
\[
 M=\|u\|_2^2,\qquad A=\|a(D)^{1/2}u\|_2^2,\qquad
 P=\operatorname{Im}\int\overline u\nabla u.
\]
We use either the local equation
\begin{equation}\label{r13g:equation}
 iu_t=a(D)u-|u|^{2m/N}u,\qquad
 p=2+2m/N,\qquad E=\frac A2-\frac{\|u\|_p^p}{p},
\end{equation}
or, when $2<m<N$, the Hartree equation
\begin{equation}\label{r13g:hartree-equation}
 iu_t=a(D)u-(K*|u|^2)u,\qquad
 E=\frac A2-\frac{H_K(u)}4,
 \quad H_K(u)=\iint K(x-y)|u(x)|^2|u(y)|^2\,dx\,dy.
\end{equation}
Here $K$ is real, even, $C^1$ away from zero, homogeneous of degree
$-m$, positive definite as a density quadratic form, and
\[
 |K(x)|+|x||\nabla K(x)|\le C_K|x|^{-m}.
\]
The corresponding sharp inequalities and mass excesses are
\begin{align*}
 \|u\|_p^p&\le\frac p2(M/M_c)^{m/N}A,
 &d&=(M/M_c)^{m/N}-1 &&\text{(local)},\\
 H_K(u)&\le2(M/M_c)A,
 &d&=M/M_c-1 &&\text{(Hartree)}.
\end{align*}
The dynamical theorem below uses strongly approximable conservative
$H^s$ solutions. No local theory for arbitrary symbols and dimensions
is silently included in that solution convention.
Positive definiteness here is strict on nonzero real densities in the
interaction space, so its square root is a strictly convex Hilbert
norm. Mere nonnegative semidefiniteness is not the assumed hypothesis.

The exterior kinetic density in both equations is
$|a(D)^{1/2}u|^2$. Convexity of $a$ is unnecessary: all anisotropic
contributions can be confined to an annulus, with an error tending to
zero. We prove that fact directly by Fourier estimates.

\subsection{An exterior kinetic term without convexity}

Put $b=\sqrt a$ and, for $\xi\ne0$, define the bounded degree-zero
symbols
\[
 B_{jk}(\xi)=\frac{\partial_j a(\xi)\xi_k}{a(\xi)}.
\]
Euler's identity gives $\sum_jB_{jj}=m$. These are real even symbols;
no sign is assigned to individual entries.

\begin{lemma}[Cutoff commutations]
\label{lem:r13g-cutoff}
For a fixed smooth compact cutoff $\theta$, $\theta_R(x)=\theta(x/R)$,
and $R\ge1$,
\begin{align}
 \|[B_{jk}(D),\theta_R]b(D)u\|_2
 &\le CR^{-1}\|u\|_{H^s},\label{r13g:zero-order-comm}\\
 \|[b(D),\theta_R]u\|_2
 &\le CR^{-1}\|u\|_{H^s}.\label{r13g:half-comm}
\end{align}
The same estimates hold for constant additions to $\theta_R$.
\end{lemma}
\begin{proof}
Write $h=\xi-\zeta$. The angular derivative bound for $B_{jk}$,
its boundedness, and $s>1$ give
\[
 |B_{jk}(\xi)-B_{jk}(\zeta)|b(\zeta)
 \le C|h||\zeta|^{s-1}.
\]
For $|h|\le|\zeta|/2$ use its homogeneous derivative of degree $-1$;
otherwise use $|\zeta|\le2|h|$. Young's convolution inequality and
$\||h|\widehat\theta_R(h)\|_1=O(R^{-1})$ prove
\eqref{r13g:zero-order-comm}. Similarly,
\[
 |b(\xi)-b(\zeta)|
 \le C|h|\bigl(|\zeta|+|h|\bigr)^{s-1}
\]
gives an upper bound
$CR^{-1}\|u\|_{H^{s-1}}+CR^{-s}\|u\|_2$ for
\eqref{r13g:half-comm}. Constants commute.
\end{proof}

Choose a radial $h_0\in C_c^\infty$ with $0\le h_0\le1$, equal
to one on $B_1$ and zero outside $B_2$. Let
\[
 f_R(x)=x h_0(x/R),\qquad
 L_R=f_R\cdot\nabla+\tfrac12\operatorname{div}f_R,
 \qquad C_R=[a(D),L_R].
\]

\begin{lemma}[Factorization of the principal form]
\label{lem:r13g-principal}
For $g=b(D)u$ and $R\ge1$,
\begin{equation}\label{r13g:factor}
 \langle u,C_Ru\rangle
 =\operatorname{Re}\sum_{j,k}
    \langle g,(\partial_j f_{R,k})B_{jk}(D)g\rangle
       +O(R^{-1}\|u\|_{H^s}^2).
\end{equation}
In particular,
\begin{align}
 \langle u,C_Ru\rangle
 &\le mA(u)-m\int_{|x|>2R}|g|^2
    +C\int_{R/2<|x|<4R}|g|^2
    +CR^{-1}\|u\|_{H^s}^2.
 \label{r13g:exterior-sign}
\end{align}
All constants are unchanged by translation of the cutoffs.
\end{lemma}
\begin{proof}
The Fourier kernel of $C_R$, with $h=\xi-\zeta$, is
\[
 \frac i2\,[a(\xi)-a(\zeta)](\xi+\zeta)_k
                                  \widehat f_{R,k}(h).
\]
The kernel of $\sum_{j,k}(\partial_j f_{R,k})q_{jk}(D)$, where
$q_{jk}(\zeta)=\partial_j a(\zeta)\zeta_k$, is
\[
 i h_j\widehat f_{R,k}(h)\partial_j a(\zeta)\zeta_k.
\]
Their difference is bounded by
\begin{equation}\label{r13g:kernel-remainder}
 C|h|^2(|\zeta|+|h|)^{m-1}|\widehat f_R(h)|.
\end{equation}
This is second-order Taylor expansion: the second derivatives of $a$
extend continuously at zero because $m>2$. To bound the form on $H^s$,
divide the kernel by $\langle\xi\rangle^s\langle\zeta\rangle^s$.
The elementary bound
\[
 \frac{(|\zeta|+|h|)^{m-1}}
      {\langle\xi\rangle^s\langle\zeta\rangle^s}
 \le C\langle h\rangle^{m-1}
\]
follows by splitting $|\zeta|\le2|h|$ and its complement.
Since $\widehat f_R(h)=R^{N+1}\widehat f_1(Rh)$, Schur's test
bounds \eqref{r13g:kernel-remainder} by $CR^{-1}$ as a form
$H^s\times H^s\to\mathbb C$.

Replacing $q_{jk}(D)$ by the $b(D)$ sandwich changes its kernel by
\[
 [b(\zeta)-b(\xi)]b(\zeta)B_{jk}(\zeta)
                      \widehat{\partial_j f_{R,k}}(h).
\]
Its magnitude is at most
$C|h|(|\zeta|+|h|)^{m-1}
|\widehat{\partial_j f_{R,k}}(h)|$.
The same weighted Schur estimate, now using the scale-invariant
coefficient $\partial_j f_{R,k}$, is $O(R^{-1})$.
Taking real parts proves \eqref{r13g:factor}.

Write
\[
 \partial_j f_{R,k}=h_0(x/R)\delta_{jk}+c_{jk,R}.
\]
The coefficients $c_{jk,R}$ are uniformly bounded and supported on
$R<|x|<2R$. The identity $\sum_jB_{jj}=m$ makes the diagonal part
exactly $m\int h_0(x/R)|g|^2$. Choose $\theta_R$ supported on
$R/2<|x|<4R$ and equal to one on the support of every $c_{jk,R}$.
For $c=c_{jk,R}$ and $B=B_{jk}(D)$,
\[
 \langle g,cBg\rangle
 =\langle\theta_Rg,cB(\theta_Rg)\rangle
   +\langle\theta_Rg,c[B,1-\theta_R]g\rangle.
\]
The first term is bounded by $C\|\theta_Rg\|_2^2$, and
Lemma~\ref{lem:r13g-cutoff} bounds the second by
$CR^{-1}\|u\|_{H^s}^2$.
Finally $h_0\le1$ and $h_0=0$ outside $2R$ give
\eqref{r13g:exterior-sign}. No pointwise principal-symbol sign has
been used.
\end{proof}

\subsection{Localized mass and moving virial identities}

\begin{lemma}[Localized mass transport]
\label{lem:r13g-mass-commutator}
For a real smooth compact cutoff $\chi_R(x)=\chi((x-z)/R)$,
$R\ge1$, and $v,w\in H^s$,
\begin{equation}\label{r13g:mass-commutator}
 |\langle v,[a(D),\chi_R]w\rangle|
 \le CR^{-1}\|v\|_{H^s}\|w\|_{H^s}.
\end{equation}
If at a given time $u=e^{i\theta}(X+iY)$, with $X,Y$ real,
$M(Y)\le CdM$, $A(Y)\le CdA(u)$, and the mass is in a fixed
compact positive interval, then for $R A(u)^{1/m}\ge1$,
\begin{equation}\label{r13g:scaled-mass-transport}
 \left|\frac d{dt}\int\chi((x-z)/R)|u|^2\right|
 \le C\sqrt d\,R^{-1}A(u)^{(m-1)/m}.
\end{equation}
The center $z$ in this statement is fixed during differentiation.
\end{lemma}
\begin{proof}
The Fourier kernel of the commutator is
$[a(\xi)-a(\zeta)]\widehat\chi_R(\xi-\zeta)$.
The mean-value bound
\[
 |a(\xi)-a(\zeta)|
 \le C|\xi-\zeta|(|\zeta|+|\xi-\zeta|)^{m-1}
\]
and the weighted Schur estimate in
Lemma~\ref{lem:r13g-principal} give \eqref{r13g:mass-commutator}.
The nonlinear potential in both equations is real and cancels from
the mass derivative. Since $a$ is real even and $\chi_R$ real,
the real-real and imaginary-imaginary terms make no contribution;
only a bilinear expression in $X,Y$ remains. Normalize $A(u)=1$
by critical scaling, apply the commutator bound and
$\|Y\|_{H^s}\le C\sqrt d$, then restore the scale.
The temporal scale contributes $A(u)$, while the cutoff radius becomes
$R A(u)^{1/m}$, giving \eqref{r13g:scaled-mass-transport}.
The chosen constant phase is never differentiated.
\end{proof}

For a locally absolutely continuous center $z(t)$, define
\[
 J_R(t)=2\operatorname{Im}\int f_R(x-z(t))\cdot\overline u\nabla u.
\]
For smooth solutions, direct differentiation gives
\begin{equation}\label{r13g:exact-moving-virial}
 J_R'=2\langle u,C_R^{z(t)}u\rangle+\mathcal N_R
 -2z'(t)\cdot\int Df_R(x-z(t))
                         \operatorname{Im}(\overline u\nabla u).
\end{equation}
The vector field $f_R=xh_0(x/R)$ has symmetric derivative, so this
matrix convention is unambiguous. The nonlinear term is
\[
 \mathcal N_R=-\frac{2(p-2)}p
       \int\operatorname{div}f_R(x-z)|u|^p
       \quad\hbox{for the local power},
\]
and
\[
 \mathcal N_R=\iint|u(x)|^2|u(y)|^2
  [f_R(x-z)-f_R(y-z)]\cdot\nabla K(x-y)\,dx\,dy
       \quad\hbox{for Hartree}.
\]
Their full-dilation values combine with $2mA$ to give $4mE$.
The motion term equals $-2z'\cdot P$ plus an error bounded by
\begin{equation}\label{r13g:motion-defect}
 C|z'|\int_{|x-z|>R}|u||\nabla u|,
\end{equation}
because $Df_R=I$ inside $B_R$. Also, for mass in a fixed compact
interval,
\begin{equation}\label{r13g:virial-endpoint}
 |J_R(t)|\le CR A(u(t))^{1/m}.
\end{equation}
Indeed $\|f_R\|_\infty\le CR$ and interpolation bounds
$\|u\|_2\|\nabla u\|_2$ by $CM^{1-1/m}A^{1/m}$.
All identities are understood in integrated form for the stated
strongly approximable solution class; the Fourier form estimates and
the interaction estimates below justify passage to the limit.

\subsection{Exterior interaction and current estimates}
\label{subsec:r13g-exterior-interaction}

The preceding estimates were stated on a bounded $H^s$ ball. Their
homogeneity gives precise instantaneous kinetic dependence. Fix the mass
in a compact positive interval, put $\lambda=A(u)^{1/m}$, and assume
$R\lambda\ge1$. Rescaling to $A=1$ in
Lemma~\ref{lem:r13g-principal} replaces its remainder by
\begin{equation}\label{r13g:scaled-principal-error}
 CR^{-1}A(u)^{(m-1)/m}.
\end{equation}
For $\zeta_R=0$ on $B_{R/4}$ and $\zeta_R=1$ outside $B_{R/2}$,
Lemma~\ref{lem:r13g-cutoff} similarly gives
\begin{equation}\label{r13g:scaled-cutoff-error}
 \|[a(D)^{1/2},\zeta_R]u\|_2^2
 \le CR^{-2}A(u)^{(m-2)/m}.
\end{equation}
These follow by unitary critical spatial rescaling; the constants do
not require an upper bound on $A(u)$.

Fix a center $z$ and write $g=a(D)^{1/2}u$,
\[
 K_R=\int_{|x-z|>8R}|g|^2,\qquad
 H_R=\int_{R/4<|x-z|<8R}|g|^2,\qquad
 \varepsilon_R=R^{-2}A(u)^{(m-2)/m}.
\]
If $\int_{|x-z|>R/4}|u|^2\le\mu$, then
\begin{equation}\label{r13g:tail-localization}
 M(\zeta_Ru)\le\mu,\qquad
 A(\zeta_Ru)\le C(K_R+H_R+\varepsilon_R).
\end{equation}
Sharp interpolation and interpolation between orders zero and $s$ give
\begin{align}
 \int_{|x-z|>R}|u|^p
 &\le C\mu^{m/N}(K_R+H_R+\varepsilon_R),
 \label{r13g:local-tail-struct}\\
 \int_{|x-z|>R}|u||\nabla u|
 &\le C\mu^{1-1/m}(K_R+H_R+\varepsilon_R)^{1/m}.
 \label{r13g:current-tail-struct}
\end{align}
For the current bound, apply
$\|\nabla w\|_2\le C\|w\|_2^{1-2/m}A(w)^{1/m}$ to
$w=\zeta_Ru$, which equals $u$ on the indicated exterior region.

\begin{lemma}[Hartree virial tail]
\label{lem:r13g-hartree-tail}
In the Hartree setting the difference between the nonlinear localized
virial expression and its full dilation value $-mH_K(u)$ obeys
\begin{equation}\label{r13g:hartree-tail}
 |\text{Hartree defect}|
 \le CR^{-m}M^2+C\mu(K_R+H_R+\varepsilon_R).
\end{equation}
\end{lemma}
\begin{proof}
For $\rho=|u|^2$ the symmetrized nonlinear expression is
\[
 \iint\rho(x)\rho(y)
 [f_R(x-z)-f_R(y-z)]\cdot\nabla K(x-y)\,dx\,dy.
\]
The full dilation vector gives $-mH_K(u)$ by homogeneity.
The difference vanishes if both points lie in $B(z,R)$, and the
uniform Lipschitz bound for $f_R$ bounds its absolute integrand by
$C\rho(x)\rho(y)|x-y|^{-m}$. By symmetry take $|x-z|>R$.
If $|y-z|<R/2$, then $|x-y|\ge R/2$, and this contribution is at
most $CR^{-m}M^2$. Otherwise both points are outside $B(z,R/2)$.
For $w=\zeta_Ru$, HLS and elliptic interpolation give
\[
 \iint|x-y|^{-m}|w(x)|^2|w(y)|^2\,dx\,dy
 \le C\|w\|_{4N/(2N-m)}^4\le CM(w)A(w).
\]
Use \eqref{r13g:tail-localization}. Thus the near-diagonal exterior
interaction sees only the small mass $\mu$, while the separated
core--tail interaction vanishes with $R$.
\end{proof}

\begin{lemma}[Hartree weak energy splitting]
\label{lem:r13g-hartree-splitting}
Suppose $u_n\rightharpoonup v$ in $H^s$ and $u_n\to v$ in
$L^2_{\mathrm{loc}}$. For $w_n=u_n-v$,
\[
 H_K(u_n)=H_K(v)+H_K(w_n)+o(1),\qquad
 E(u_n)=E(v)+E(w_n)+o(1).
\]
\end{lemma}
\begin{proof}
Set $q=4N/(2N-m)$, strictly between $2$ and the $H^s$ Sobolev
exponent. Thus $w_n\to0$ in $L^q_{\mathrm{loc}}$ and is bounded
globally in $L^q$. Approximate $v$ by smooth compact functions to
obtain $\operatorname{Re}(\overline v w_n)\to0$ in $L^{q/2}$.
HLS removes all corresponding terms in the density expansion.
It remains to show
$\iint K(x-y)|v(x)|^2|w_n(y)|^2\,dx\,dy\to0$.
The potential of a smooth compact approximation to $v$ is bounded
because $m<N$, and tends to zero at infinity by the Riesz bound.
Local strong $L^2$ convergence and bounded mass of $w_n$ make this
pairing tend to zero. HLS controls the approximation error.
The kinetic quadratic form splits by weak convergence, proving the
energy identity.
\end{proof}

\subsection{A concrete conservative local flow}

\begin{proposition}[Integer energy regularity in dimensions one to three]
\label{prop:r13g-local-flow}
Let $m=2k$ with integer $k\ge2$, $N=1,2,3$, and let $\mathcal H$
be any separable complex Hilbert space. For every symbol $a$ above,
let $b$ be a common real measurable symbol satisfying
$0\le b(\xi)\le C_b(1+|\xi|^{2k})$. The total-density system
\[
 i\partial_tU=[a(D)+b(D)]U-\|U\|_{\mathcal H}^{4k/N}U
\]
has a unique maximal $C_tH^k_x(\mathcal H)$ solution for every
$H^k(\mathcal H)$ datum. The flow conserves
\begin{gather*}
 M=\int\|U\|_{\mathcal H}^2,\qquad
 P=\operatorname{Im}\int\langle U,\nabla U\rangle_{\mathcal H},\\
 E_b=\frac12\int(a+b)(\xi)\|\widehat U(\xi)\|_{\mathcal H}^2\,d\xi
       -\frac1{\alpha+2}\int\|U\|_{\mathcal H}^{\alpha+2},
 \qquad\alpha=4k/N,
\end{gather*}
and has the usual $H^k$ continuation alternative. The local estimates
are independent of the internal dimension. For the homogeneous case
$b=0$, the localized mass and virial identities also pass by strong
approximation. In particular the scalar dynamical conclusions below
need no additional local-theory assumption throughout this integer-order
class, which includes every positive elliptic even homogeneous
polynomial of degree $2k$. No homogeneous virial identity is asserted
for a general nonzero $b$.
\end{proposition}
\begin{proof}
Write $\alpha=4k/N>k$. On the underlying real Hilbert space, the
radial map $F(Z)=\|Z\|_{\mathcal H}^{\alpha}Z$ is Fr\'echet
$C^{k+1}$. Its derivatives obey, for $1\le j\le k+1$,
\[
 \|D^jF(Z)\|_{\mathrm{op}}
 \le C_{\alpha,j}\|Z\|_{\mathcal H}^{\alpha+1-j},
\]
with constants independent of $\mathcal H$; the derivatives extend
continuously at zero because $\alpha+1-j>0$. These bounds follow
by differentiating the Hilbert norm away from zero: each term is a
product of inner products, vectors, and a power of the norm.
Since $k>N/2$, the Hilbert-valued Sobolev embedding and integer
chain/product estimates, applied to
\[
 F(U)-F(V)=\int_0^1 DF(V+\tau(U-V))(U-V)\,d\tau,
\]
show that $F:H^k(\mathcal H)\to H^k(\mathcal H)$ is locally
Lipschitz. The estimates use operator norms of the displayed
derivatives and pointwise Hilbert norms, so their constants do not
grow with the internal dimension. They extend from finite-component
smooth fields by orthogonal projection and density. The common scalar
group $e^{-it[a(D)+b(D)]}$ is strongly continuous and unitary on
$H^k(\mathcal H)$. Duhamel contraction gives existence, uniqueness,
continuous dependence, and the blowup alternative.

For conservation and approximation, use the self-adjoint Fourier
projection $\Pi_n$ onto $|\xi|\le n$ and solve
\[
 i\partial_tU_n=[a(D)+b(D)]U_n-\Pi_nF(U_n),\qquad
 U_n(0)=\Pi_nU_0,\qquad U_n=\Pi_nU_n.
\]
This projected Hamiltonian flow conserves the restrictions of
$M,P,E_b$: $\Pi_n$ commutes with translations and the linear
dispersion, and its range is invariant under constant phase.
The local Lipschitz bounds are uniform in $n$, so Duhamel and
Gronwall give $U_n\to U$ in $C_tH^k_x(\mathcal H)$ on every compact subinterval
of the lifespan. All three conserved functionals are continuous in
$H^k(\mathcal H)$, using the stated growth bound for $b$, so
conservation passes to the limit. When $b=0$, the same approximation
proves the integrated localized identities: the linear commutators
are continuous $H^k$ forms by the Fourier bounds above, which also
hold componentwise with dimension-independent constants. The nonlinear
projection error vanishes because $F(U_n)\to F(U)$ in $H^k$ and
the fixed-radius first-order virial operator is bounded from $H^k$
to $L^2$. No dispersive decay estimate is required.
\end{proof}


\section{Persistent amplification for homogeneous dispersion}
\label{sec:r13-persistent-general}

The exterior commutator estimates yield a stronger statement than one
kinetic excursion: a lower long-time average in actual physical time,
without a global kinetic bound. Let $m>2$, $s=m/2$, and let $a$ satisfy
the positive elliptic, real-even, homogeneous symbol assumptions above.
Write $A(u)=\|a(D)^{1/2}u\|_2^2$, $M(u)=\|u\|_2^2$ and
$P(u)=\operatorname{Im}\int\overline u\nabla u$.

We treat both the local equation $iu_t=a(D)u-|u|^{2m/N}u$ and the
Hartree equation $iu_t=a(D)u-(K*|u|^2)u$ under the preceding Hartree
kernel assumptions: $2<m<N$, a real-even, degree-$-m$ homogeneous
$C^1$ kernel away from zero, positive-definite density form, and
$|K(x)|+|x||\nabla K(x)|\le C|x|^{-m}$. Put
\[
 \kappa=1+m/N\quad\hbox{(local)},\qquad
 \kappa=2\quad\hbox{(Hartree)}.
\]
Their sharp energy inequalities have the common form
\[
 E(u)=A(u)/2-\mathcal F(u)/(2\kappa),\qquad
 \mathcal F(u)\le\kappa(M(u)/M_c)^{\kappa-1}A(u).
\]
Here $\mathcal F=\|u\|_{2\kappa}^{2\kappa}$ locally and is the Hartree
interaction otherwise.

\begin{theorem}[Lower long-time kinetic average]\label{thm:r13-general-average}
For each fixed equation above there are $d_0,c>0$ such that every
forward-global conservative strong $H^s$ solution, for which the
localized identities pass by strong energy-space approximation, with
\[
 P=0,\qquad E=-e<0,\qquad
 0<d=(M/M_c)^{\kappa-1}-1<d_0
\]
satisfies
\begin{equation}\label{eq:r13-general-average}
 \boxed{\quad
 \liminf_{T\to\infty}\frac1T\int_0^T A(u(t))\,dt
 \ge c e d^{-\alpha_m},\qquad
 \alpha_m=1+\frac{m}{2(m-1)}.
 \quad}
\end{equation}
No uniform kinetic bound, radial symmetry, finite spatial moment,
convexity of $a$, or real-sector linearized nondegeneracy is assumed.
The theorem specifies its solution class; it does not assert local
well-posedness for every symbol and dimension.
\end{theorem}

\begin{proof}
Write $A(t)=A(u(t))$, $\delta=M-M_c$, $\gamma=(m-1)/m$.
Sharp interpolation gives $A(t)\ge a_*:=2e/d>0$.
Real polarization supplies, separately at each time, a phase with
$u=e^{i\theta}(X+iY)$, $X,Y$ real, and
$M(Y)\le CdM$, $A(Y)\le CdA(t)$.
The nonlinear potential is real in both models and cancels from the
localized mass derivative. Apply the normalized Fourier mass-commutator
bound after scaling $A$ to one. For $R A(t)^{1/m}\ge1$ it gives
\begin{equation}\label{eq:r13g-variable-density}
 \left|\frac d{dt}\int\chi((x-z)/R)|u|^2\right|
 \le C\sqrt d R^{-1}A(t)^\gamma.
\end{equation}
No time derivative of the chosen phase is taken.

\medskip\noindent\emph{Center and uniform mass tail.}
Fix $\varepsilon_0=M_c/32$ and take $d_0$ small. Normalized
near-optimizer compactness gives a fixed $\overline R$ and centers
$c(t)$ with
\[
 \int_{|x-c(t)|\le\overline R A(t)^{-1/m}}|u(t,x)|^2\,dx
 \ge M_c-\varepsilon_0.
\]
Choose $R_0\ge1$ containing these balls and satisfying
$R_0\ge a_*^{-1/m}$. Introduce $s(t)=\int_0^tA(\tau)^\gamma d\tau$.
Two centers separated by more than $4R_0$ force a cutoff mass change
of at least $M_c-3\varepsilon_0$. Equation
\eqref{eq:r13g-variable-density} excludes this when
$|s(t)-s(t')|\le\tau R_0/\sqrt d$, with fixed small $\tau>0$.
Interpolating sampled centers at these intrinsic-time intervals gives
a locally absolutely continuous $b(t)$ with
\begin{equation}\label{eq:r13g-variable-center}
 |b(t)-c(t)|\le8R_0,\qquad
 |b'(t)|\le C\sqrt d A(t)^\gamma\quad\hbox{a.e.}
\end{equation}
The speed constant is independent of $R_0$ and $e/d$.

For every $\eta>0$ there is $R_\eta<\infty$ such that
\begin{equation}\label{eq:r13g-variable-tail}
 \sup_{t\ge0}\int_{|x-b(t)|>R_\eta}|u(t,x)|^2\,dx
 \le\delta+\eta.
\end{equation}
Here is the compactness argument without a kinetic upper bound.
Suppose bad times $t_k$ and radii $R_k\to\infty$ exist. The bad mass
lies also outside $B(c(t_k),R_k-8R_0)$. Set $A_k=A(t_k)$,
$\ell_k=A_k^{-1/m}$ and
$v_k(y)=\ell_k^{N/2}u(t_k,c(t_k)+\ell_k y)$.
Then $A(v_k)=1$, $M(v_k)=M$, $E(v_k)=-e/A_k\in[-d/2,0)$,
and a fixed ball contains $M_c-\varepsilon_0$ mass.
A weak $H^s$ limit $v$, with strong local $L^2$ convergence, therefore
has $M(v)\ge M_c-\varepsilon_0>0$. The normalized bad radius tends
to infinity, since
$(R_k-8R_0)\ell_k^{-1}\ge(R_k-8R_0)a_*^{1/m}\to\infty$.
Consequently $M(v)\le M_c-\eta$.
For $w_k=v_k-v$, mass and energy split:
\[
 E(v_k)=E(v)+E(w_k)+o(1),\qquad
 M(w_k)\le\delta+\varepsilon_0+o(1)<M_c.
\]
Locally the nonlinear splitting is Brezis--Lieb. For Hartree, the
Riesz bound and HLS imply that $v\overline{w_k}$ vanishes in the
density space, while $|w_k|^2$ converges weakly to zero there. First
approximate $v$ by bounded compactly supported functions and use local
compactness; then use uniform Sobolev and HLS bounds. This removes all
mixed interaction terms. Sharp interpolation gives $E(w_k)\ge0$
eventually and $E(v)>0$, because $v\ne0$ is strictly subcritical,
contradicting $E(v_k)<0$. This proves \eqref{eq:r13g-variable-tail}.

\medskip\noindent\emph{Variable-scale virial.}
Let $g=a(D)^{1/2}u$ and use the compact vector field $f_R$ of the
exterior commutator lemma. Define
\[
 J_R=2\operatorname{Im}\int f_R(x-b(t))\cdot\overline u\nabla u,
 \quad K_R=\int_{|x-b|>8R}|g|^2,
 \quad H_R=\int_{R/4<|x-b|<8R}|g|^2.
\]
Scaling the normalized exterior commutator and cutoff-energy estimates
back to $A(t)$ gives errors $CR^{-1}A^\gamma$ and
$CR^{-2}A^{(m-2)/m}$. Write
\[
 q_R(t)=R^{-1}A(t)^\gamma+R^{-2}A(t)^{(m-2)/m}+R^{-m}.
\]
The last term is needed only for separated Hartree interactions.
For $R\ge4R_\eta$, $R\ge a_*^{-1/m}$ and $\mu=\delta+\eta$,
the preceding localized identities give
\begin{align*}
 J_R'\le{}&-4me-2mK_R+CH_R+Cq_R\\
 &+C\mu^{\kappa-1}(K_R+H_R+q_R)\\
 &+C|b'|\mu^{1-1/m}(K_R+H_R+q_R)^{1/m}.
\end{align*}
The constant center-motion term is $-2b'\cdot P=0$. Its remainder is
bounded using $\|\zeta_Ru\|_2\|\nabla(\zeta_Ru)\|_2$ and Fourier
interpolation, since $s>1$. For Hartree, separated inner--outer
interactions are $O(R^{-m}M^2)$; the remaining exterior interaction
is bounded by $C\mu A(\zeta_Ru)$.

Take $d_0,\eta$ small enough to absorb the nonlinear $K_R$ term.
Young's inequality absorbs the current's $K_R$ term and bounds the
other parts by $C(H_R+q_R)+C\mu|b'|^{m/(m-1)}$.
Equation \eqref{eq:r13g-variable-center} implies
\begin{equation}\label{eq:r13g-average-virial}
 J_R'\le-4me+CH_R+C\mu d^{m/[2(m-1)]}A(t)+Cq_R(t).
\end{equation}
The endpoint estimate is $|J_R(t)|\le CR A(t)^{1/m}$.

\medskip\noindent\emph{Physical-time averaging.}
Let $L$ be the left side of \eqref{eq:r13-general-average}; only finite
$L$ needs proof. Choose $T_n\to\infty$ with averages tending to $L$.
There is $t_n\in[T_n-\sqrt{T_n},T_n]$ with
$A(t_n)\le C(L+1)\sqrt{T_n}$. Therefore $t_n/T_n\to1$, the
averages through $t_n$ tend to $L$, and $A(t_n)^{1/m}/t_n\to0$.
For fixed $K<\infty$ and large $R_*$, the annuli at radii
$R_j=64^jR_*$, $0\le j<K$, are disjoint, so
$\sum_jH_{R_j}(t)\le A(t)$.
For each $n$ choose one radius with average annular energy at most
$K^{-1}$ times the average of $A$ and integrate
\eqref{eq:r13g-average-virial}.
The chosen radii stay bounded for fixed $K,R_*$, so endpoints divided
by $t_n$ vanish. Jensen's inequality bounds the fractional powers of
$A$ in $q_R$. The time limit yields
\[
 4me\le CL/K+C\mu d^{m/[2(m-1)]}L
       +CR_*^{-1}L^\gamma+CR_*^{-2}L^{(m-2)/m}+CR_*^{-m}.
\]
Let $K\to\infty$, then $R_*\to\infty$, then $\eta\downarrow0$.
Since $\delta=M_c[(1+d)^{1/(\kappa-1)}-1]\le Cd$, this proves
\eqref{eq:r13-general-average}.
\end{proof}

\begin{corollary}[Repeated amplification and standing-wave instability]
\label{cor:r13-general-repeat}
For any real sharp standing wave $Q$ with $M(Q)=M_c$, take
$u_d(0)=(1+d)^{1/[2(\kappa-1)]}Q$.
Every global solution in the stated class satisfies
\[
 \liminf_{T\to\infty}\frac1T\int_0^T
  \frac{A(u_d(t))}{A(u_d(0))}\,dt
 \ge c d^{-m/[2(m-1)]}.
\]
Every prescribed finite amplification factor is thus exceeded at
arbitrarily late times for small enough $d$. Under the usual local
continuation alternative, each real sharp standing wave is nonlinearly
orbitally unstable modulo phase and translation.
\end{corollary}
\begin{proof}
The exact amplitude identity gives $e=dA(u_d(0))/2$.
The initial perturbations converge to $Q$ in $H^s$, whereas a large
kinetic excursion leaves a fixed neighborhood of its phase--translation
orbit. Finite-time breakdown also forces such an excursion under the
continuation alternative.
\end{proof}

If a particular orbit has $B=\sup_tA(t)<\infty$ and
$L_*=ce d^{-\alpha_m}$, then the set $\{t:A(t)\ge L_*/2\}$ has
lower asymptotic density at least $L_*/(2B)$. Without this extra
boundedness, the theorem still permits sparse high spikes. For $m=4$,
$\alpha_m=5/3$ and the perturbation exponent is $2/3$. The result does
not by itself exclude every bounded global negative-energy orbit.

\section{Momentum-adapted centers and stronger amplification}
\label{sec:r13-center}

The preceding amplification theorem requires no nondegeneracy of the
real linearized operator. We now show what that additional hypothesis
buys dynamically: conserved momentum removes the leading translation
speed in a suitable modulation coordinate. The scale and phase
equations must be retained for this cancellation to be useful.

Throughout this section $N\in\{1,2,3\}$ and
\[
 iu_t=\Delta^2u-|u|^{8/N}u,\qquad p=2+8/N.
\]
Let $Q$ be a real sharp ground wave satisfying
\[
 \Delta^2Q+Q-|Q|^{p-2}Q=0,\qquad M(Q)=M_c,
\]
and put
\[
 L_-=\Delta^2+1-|Q|^{p-2},\quad
 L_+=\Delta^2+1-(p-1)|Q|^{p-2},\quad
 \Lambda=N/2+x\cdot\nabla.
\]
We assume
\begin{equation}\label{eq:r13-real-nondeg}
 \ker L_+=\operatorname{span}\{\partial_jQ:1\le j\le N\}.
\end{equation}
The phase coercivity and simplicity of $\ker L_-$ are provided by the
sharp variational problem. The weighted regularity of the profile and
dual functions needed below follows from this concrete elliptic
equation, as we prove next; it is not an additional assumption.

\begin{lemma}[Exponential regularity of the dual functions]
\label{lem:r13-dual-decay}
For some $a>0$, $e^{a\langle x\rangle}Q\in H^6$. If
$L_-Z=f$ on $Q^\perp$, with $f=\partial_jQ$ or $f=\Lambda Q$,
then $e^{a'\langle x\rangle}Z\in H^6$ for some $a'>0$.
In particular $\Lambda Q\in H^4$, and the functions $Z_j,W$ below
satisfy $\nabla Z_j,\Lambda Z_j,\nabla W,\Lambda W\in L^2$.
\end{lemma}
\begin{proof}
The $H^2$ algebra property and the profile equation give $Q\in H^6$.
Hence $V=|Q|^{8/N}\in W^{2,\infty}$ and $V$ tends to zero at infinity.
For smooth bounded truncations $\phi_R$ of $a\langle x\rangle$, put
$w_R=e^{\phi_R}Q$. The exact identity
\[
 \operatorname{Re}\langle\Delta^2Q,e^{2\phi_R}Q\rangle
 =\|\Delta w_R+|\nabla\phi_R|^2w_R\|_2^2
 -\|2\nabla\phi_R\cdot\nabla w_R+(\Delta\phi_R)w_R\|_2^2
\]
differs from $\|\Delta w_R\|_2^2$ by at most
$Ca^2\|w_R\|_{H^2}^2$. The potential is small outside a fixed ball;
its weighted integral inside that ball is uniformly bounded.
Absorption therefore gives a uniform weighted $H^2$ bound. The
conjugated equation has smooth bounded lower-order coefficients of
order at most three. Since multiplication by $V$ preserves $H^2$,
successive elliptic gains give
$e^{a\langle x\rangle}Q\in H^6$.

For $L_-Z=f$, choose $0<a'<a$ small enough and truncate
$a'\langle x\rangle$ with $\phi_R\le a'\langle x\rangle$.
The same test produces the source term
\[
 |\langle e^{\phi_R}f,e^{\phi_R}Z\rangle|
 \le \tfrac14\|e^{\phi_R}Z\|_2^2
      +C\|e^{a'\langle x\rangle}f\|_2^2.
\]
The rightmost norm is finite, and the interior weighted potential is
bounded by the fixed unweighted $L^2$ norm of $Z$ on the ball.
Absorption and passage to the limit again give weighted $H^2$.
Moreover $e^{a'\langle x\rangle}f\in H^2$. The conjugated equation
then gives the successive gains $H^2\to H^3\to H^4\to H^5\to H^6$.
Polynomially weighted derivatives of lower order follow by reducing
the exponential weight. Differentiating the scaled profile equation
now also justifies $L_+\Lambda Q=-4Q$.
\end{proof}

\begin{theorem}[Quadratic drift and stronger amplification]
\label{thm:r13-quadratic-drift}
Under these hypotheses there are constants $d_Q,C_Q>0$ and a
neighborhood $\mathcal U_Q$ of the normalized phase--translation orbit
of $Q$ such that the following holds. Let $u$ be a conservative,
strongly approximable $H^2$ solution with the usual continuation
property, and suppose
\[
 0<d=\left(\frac{M}{M_c}\right)^{4/N}-1<d_Q,
 \qquad E=-e<0,\qquad P=0.
\]
Assume the datum, after amplitude normalization to mass $M_c$ and
spatial normalization to kinetic energy $A(Q)$, lies in
$\mathcal U_Q$. Every forward-global such solution satisfies
\begin{equation}\label{eq:r13-seven-thirds-average}
 \liminf_{T\to\infty}\frac1T\int_0^T A(u(t))\,dt
 \ge c_Q e d^{-7/3}.
\end{equation}
In particular, if $B=\sup_{t\ge0}A(u(t))<\infty$, then
\begin{equation}\label{eq:r13-seven-thirds}
 e\le C_Q d^{7/3}B.
\end{equation}
In particular, for
$u_d(0)=(1+d)^{N/8}Q$,
\begin{equation}\label{eq:r13-four-thirds}
 \sup_{0\le t<T_+}\frac{A(u_d(t))}{A(u_d(0))}
 \ge c_Qd^{-4/3}.
\end{equation}
Every factor strictly smaller than the right-hand side is attained at
a finite time, and at arbitrarily late times if the solution is global.
In that case the lower long-time average of $A(u_d(t))/A(u_d(0))$ is
also at least $c_Qd^{-4/3}$.
No radial condition is imposed on the evolving solution.
\end{theorem}

\begin{proof}
Define
\[
 Z_j=L_-^{-1}\partial_jQ,\quad
 T_{ij}=\langle\partial_iQ,Z_j\rangle,\quad
 c_j=\langle\Lambda Q,Z_j\rangle.
\]
Phase coercivity implies $T>0$. Translating $Q(x)$ to $Q(x+b)$ changes
$c$ to $c-Tb$. Choose $b=T^{-1}c$, so that $c=0$, and define
\[
 W=L_-^{-1}\Lambda Q,\qquad D=\langle\Lambda Q,W\rangle>0.
\]
Selfadjointness gives
\begin{equation}\label{eq:r13-dual-orth}
 \langle\Lambda Q,Z_j\rangle=
 \langle\partial_jQ,W\rangle=0.
\end{equation}
Strict positivity of $D$ follows from phase coercivity and
$\Lambda Q\ne0$; a nonzero $L^2$ function cannot be homogeneous of
degree $-N/2$.

\emph{Coercivity and isolation.}
Sharp interpolation implies $E(f)\ge0$ on $M(f)=M_c$. Therefore
$\langle L_+h,h\rangle\ge0$ on $Q^\perp$. A zero vector of this
constrained form satisfies $L_+h\in\mathbb RQ$. Since
$L_+\Lambda Q=-4Q$, assumption \eqref{eq:r13-real-nondeg} identifies
the constrained null space as
$\operatorname{span}\{\Lambda Q,\partial_1Q,\ldots,\partial_NQ\}$.
The decaying potential is form-compact relative to $\Delta^2+1$.
Consequently, \eqref{eq:r13-dual-orth} and $T>0$ yield
\begin{equation}\label{eq:r13-constrained-coercivity}
 \langle L_+h,h\rangle+\langle L_-k,k\rangle
 \ge c_Q(\|h\|_{H^2}^2+\|k\|_{H^2}^2)
       -C_Q\langle h,Q\rangle^2
\end{equation}
whenever
\begin{equation}\label{eq:r13-gauges}
 \langle h,Z_j\rangle=\langle h,W\rangle=
 \langle k,Q\rangle=0.
\end{equation}
The same estimate, applied to a nearby exact optimizer with its exact
mass and energy constraints, shows that the normalized orbit of $Q$
is isolated. Indeed its perturbation norm is bounded by a higher-order
remainder in that same norm and must vanish.

For a negative-energy trajectory of mass parameter $d$, the normalized
sharp-interpolation deficit is at most $d$. Sharp-optimizer compactness
therefore places the normalized trajectory uniformly close to the set
of normalized optimizers as $d\to0$. Choose $\rho>0$ so that the
quotient ball of radius $2\rho$ about $Q$'s orbit contains no other
normalized optimizer orbit. A state at distance $\rho$ from that orbit
is at distance at least $\rho$ from every optimizer orbit, by the
triangle inequality. This contradicts uniform distance less than
$\rho/4$ from the optimizer set for small $d$. Thus a continuous
normalized trajectory starting at distance less than $\rho/2$ cannot
leave the $\rho$ tube. Choosing $\rho$ small enough proves that one
fixed $Q$ can be used throughout the maximal existence interval.

\emph{Exact modulation equations.}
Write
\[
 u(t,x)=\lambda(t)^{N/2}e^{i\theta(t)}w(s,y),\qquad
 w=Q+h+ik,\qquad y=\lambda(t)(x-a(t)),\qquad
 \frac{ds}{dt}=\lambda^4,
\]
with \eqref{eq:r13-gauges}. The parameter Jacobian at $Q$ has
invertible blocks $T,D,M_c$. Lemma~\ref{lem:r13-dual-decay} justifies
local absolute continuity of the parameters and
differentiation of the gauges. Terms involving $\Lambda h$ are paired
by integration by parts against $\Lambda Z_j$ or $\Lambda W$.
Set
\[
 b=\lambda_s/\lambda,\qquad v=\lambda a_s,\qquad
 \alpha=\theta_s-1,
\]
and, for $F(z)=|z|^{p-2}z$, set
\[
 \mathcal R=F(Q+h+ik)-F(Q)
 -(p-1)|Q|^{p-2}h-i|Q|^{p-2}k.
\]
The exact real equations are
\begin{align}
 h_s&=L_-k-\operatorname{Im}\mathcal R+\alpha k
       -b\Lambda(Q+h)+v\cdot\nabla(Q+h),\label{eq:r13-real-flow}\\
 k_s&=-L_+h+\operatorname{Re}\mathcal R-\alpha(Q+h)
       -b\Lambda k+v\cdot\nabla k.\label{eq:r13-imag-flow}
\end{align}
If $\epsilon^2=\|h\|_{H^2}^2+\|k\|_{H^2}^2$, then
$\|\mathcal R\|_2\le C_Q\epsilon^2$, since $N\le3$, $p>4$, and
$H^2\hookrightarrow L^\infty$.

Zero momentum is exactly
\begin{equation}\label{eq:r13-momentum-cancel}
 \langle k,\partial_jQ\rangle=-\langle k,\partial_jh\rangle.
\end{equation}
Let $(T_h)_{j\ell}=\langle\partial_\ell h,Z_j\rangle$.
Differentiate the translation, scale and phase gauges in that order:
\begin{align}
 [(T+T_h)v]_j
 &=\langle k,\partial_jh\rangle
   +\langle\operatorname{Im}\mathcal R,Z_j\rangle
   -\alpha\langle k,Z_j\rangle
   +b\langle\Lambda h,Z_j\rangle,
   \label{eq:r13-translation-eq}\\
 b[D+\langle\Lambda h,W\rangle]-v\cdot\langle\nabla h,W\rangle
 &=\langle k,\Lambda Q\rangle
   -\langle\operatorname{Im}\mathcal R,W\rangle
   +\alpha\langle k,W\rangle,
   \label{eq:r13-scale-eq}\\
 \alpha[M_c+\langle h,Q\rangle]+b\langle\Lambda k,Q\rangle
   -v\cdot\langle\nabla k,Q\rangle
 &=-\langle h,L_+Q\rangle
   +\langle\operatorname{Re}\mathcal R,Q\rangle.
   \label{eq:r13-phase-eq}
\end{align}
This linear system for $v,b,\alpha$ is a small perturbation of one
with diagonal blocks $T,D,M_c$, so
$|v|+|b|+|\alpha|\le C_Q\epsilon$. Substitution in
\eqref{eq:r13-translation-eq} improves the translation estimate to
\begin{equation}\label{eq:r13-quadratic-v}
 |v|\le C_Q\epsilon^2.
\end{equation}
The phase and scale terms have thus been controlled, not omitted.

\emph{Energy and center speed.}
Put $\delta M=M-M_c=O(d)$. The mass identity gives
$2\langle h,Q\rangle=\delta M-\|h\|_2^2-\|k\|_2^2$.
For $S=E+M/2$, Taylor expansion and
\eqref{eq:r13-constrained-coercivity} imply
\[
 S(w)-S(Q)\ge c_Q\epsilon^2-C_Q(\delta M)^2-o(\epsilon^2).
\]
But $E(w)=\lambda^{-4}E(u)<0$, so
$S(w)-S(Q)\le\delta M/2$. Shrinking the tube yields
\begin{equation}\label{eq:r13-small-error}
 \epsilon^2\le C_Qd,\qquad
 |a_t|=\lambda^3|v|\le C_Qd\lambda^3.
\end{equation}

Since $A(w)$ stays comparable to $A(Q)$,
\eqref{eq:r13-small-error} gives the scale-sensitive speed
\[
 |a_t|\le C_Qd A(t)^{3/4}
\]
without any upper bound on $A(t)$. A ball of radius $C_Q A(t)^{-1/4}$
about $a(t)$ contains at least $M_c-\varepsilon_0$ mass, where
$\varepsilon_0>0$ is fixed and small. The bound $A(t)\ge2e/d$ makes
these physical radii uniformly bounded.

For every $\eta>0$ the mass outside a sufficiently large fixed radius
about $a(t)$ is at most $M-M_c+\eta$, uniformly in time. To justify
this without a kinetic bound, suppose bad times and radii tend to
infinity and normalize each state to kinetic energy one about its
center. The bad normalized radii tend to infinity because
$A(t)\ge2e/d$. Weak $H^2$ compactness gives a nonzero core of mass
between $M_c-\varepsilon_0$ and $M_c-\eta$. The remainder is also
subcritical. Their split energies have strictly positive sum, in
contradiction with the normalized total energies $-e/A(t)\le0$.

Put $\mu=M-M_c+\eta$. The negative exterior kinetic contribution
of Lemma~\ref{lem:r13g-principal}, with $a(\xi)=|\xi|^4$,
absorbs the exterior current by Young's
inequality. With $H_R$ denoting the transition-annulus kinetic energy,
the resulting estimate is
\[
 J_R'\le-16e+CH_R+C_Q\mu d^{4/3}A(t)
                  +C[R^{-1}A(t)^{3/4}+R^{-2}A(t)^{1/2}+R^{-4}].
\]
Let $L=\liminf_{T\to\infty}T^{-1}\int_0^T A$. If $L=\infty$ the
claim follows. Otherwise choose times $t_n\to\infty$ with the
averages tending to $L$ and $A(t_n)^{1/4}/t_n\to0$; the endpoint
selection in the preceding actual-time argument applies unchanged.
Since $|J_R(t)|\le CR A(t)^{1/4}$, endpoint terms disappear after
division by $t_n$. Among $K$ disjoint geometrically spaced annuli,
one has average kinetic energy at most the total average divided by
$K$. Integrate at that radius, then let time tend to infinity, $K$
tend to infinity, the base radius tend to infinity, and $\eta$ tend
to zero. This gives
\[
 e\le C_Q(M-M_c)d^{4/3}L\le C_Qd^{7/3}L.
\]
It proves \eqref{eq:r13-seven-thirds-average}, and
\eqref{eq:r13-seven-thirds} follows immediately.
For the specified initial data $E=-dA(u_d(0))/2$ and $P=0$.
The average bound gives arbitrarily late crossings for global
solutions. Finite-time breakdown makes the kinetic supremum infinite
by continuation. Continuity gives the stated finite-time crossings.
\end{proof}

\begin{theorem}[Arbitrary momentum in the same ground component]
\label{thm:r13-any-momentum}
Retain the hypotheses of Theorem~\ref{thm:r13-quadratic-drift}, except
that $P$ may now be arbitrary. Every forward-global such solution
satisfies
\begin{equation}\label{eq:r13-any-momentum-average}
 \liminf_{T\to\infty}\frac1T\int_0^T A(u(t))\,dt
 \ge c_Q|E|d^{-2}.
\end{equation}
At $P=0$, \eqref{eq:r13-seven-thirds-average} gives the stronger
exponent $7/3$.
\end{theorem}
\begin{proof}
Component trapping, constrained coercivity and
$\epsilon^2\le C_Qd$ did not use $P=0$. For general momentum the
exact identity replacing \eqref{eq:r13-momentum-cancel} is
\[
 \langle k,\partial_jQ\rangle
 =-\frac{P_j}{2\lambda}-\langle k,\partial_jh\rangle.
\]
The right-hand side of \eqref{eq:r13-translation-eq} therefore gains
$P_j/(2\lambda)$. Since $|P|/\lambda\le C_Q\epsilon$, solving the
same coupled equations first to linear order and substituting back
gives
\begin{equation}\label{eq:r13-effective-drift}
 v=\tfrac12T^{-1}(P/\lambda)+O_Q(d),\qquad
 a_t=\tfrac{\lambda^2}{2}T^{-1}P+O_Q(d\lambda^3).
\end{equation}
Retain the exact $-2a_t\cdot P$ term in localized virial. Positivity
of $T$ gives
\[
 -2a_t\cdot P\le-c_Q\lambda^2|P|^2+C_Qd^2\lambda^4.
\]
After absorbing exterior kinetic energy, the exterior current costs
$C\mu|a_t|^{4/3}$, where $\mu=M-M_c+\eta$. For any fixed small
$\varepsilon_1>0$, \eqref{eq:r13-effective-drift} and Young's
inequality imply
\begin{align*}
 \mu|a_t|^{4/3}
 &\le C_Q\mu\lambda^{8/3}|P|^{4/3}
       +C_Q\mu d^{4/3}\lambda^4\\
 &\le\varepsilon_1\lambda^2|P|^2
       +C_{Q,\varepsilon_1}(\mu^3+\mu d^{4/3})\lambda^4.
\end{align*}
The first term is absorbed into the negative momentum contribution.
Since $\lambda^4$ is comparable to $A(t)$, this leaves
\[
 J_R'\le-16e+CH_R
   +C_Q[d^2+\mu^3+\mu d^{4/3}]A(t)
   +C[R^{-1}A(t)^{3/4}+R^{-2}A(t)^{1/2}+R^{-4}].
\]
Uniform surplus-tail mass and actual-time endpoint selection remain
valid. Averaging disjoint annuli and sending $\eta\to0$ gives
\[
 e\le C_Q[d^2+(M-M_c)^3+(M-M_c)d^{4/3}]
       \liminf_{T\to\infty}\frac1T\int_0^T A(t)\,dt.
\]
Because $M-M_c=O(d)$, the bracket is $O(d^2)$. This proves
\eqref{eq:r13-any-momentum-average}. No Galilean invariance is used.
\end{proof}

The conditional exponent $4/3$ improves the unconditional exponent
$2/3$. It still permits a bounded trajectory at a sufficiently large
kinetic scale. Reapplying the estimate later leaves $d$ and $E$
unchanged and does not force iteration to infinity. A further dynamical
argument must control scale evolution or the quadratic translation
terms in \eqref{eq:r13-translation-eq} over long times.


\section{Relation to existing work and the remaining evolution problem}
For the pure quartic equation, Boulenger--Lenzmann~\cite{BL} prove
radial negative-energy finite- or infinite-time blowup, and
Bonheure--Cast\'eras--Gou--Jeanjean~\cite{BCGJ} deduce ground-wave
instability under their stated radial hypotheses. The recent
Dinh--Muse manuscript~\cite{DM}, Theorem 1.1 and Remark 1.4, states
radial finite-time blowup and leaves the unrestricted nonradial
question open. Instability in the full perturbation space is already
implied by instability using radial perturbations of the same wave.
Thus orbital instability alone should not be presented as resolving
the nonradial blowup problem. The distinctive conclusion studied here
is a quantitative physical-time average for general near-threshold
trajectories, with the precise momentum and spectral assumptions
stated in each theorem.

Pausader--Shao~\cite{PS}, Section 6, already relate center motion,
momentum, and energy for fixed-scale precompact quartic trajectories
with the additional regularity in their proposition. Their rigidity
argument rules out the corresponding negative-energy zero-momentum
compact trajectory. Our tail argument permits small mass radiation
and does not assume kinetic compactness. The remaining moving-tail
term is exactly what prevents the present theorem from giving full
compact-trajectory rigidity for every bounded solution.

Inui--Komada--Nishimura~\cite{IKN} prove nonradial subground scattering
in a range starting at nonlinearity power $1+8/(N-1)$, using a
virial transverse to momentum. The mass-critical power $1+8/N$ is
outside that range: its transverse scaling exponent is below four.
Their argument therefore does not directly settle the critical
problem addressed here.

Occupation-sensitive interpolation has a substantial history,
including Dolbeault--Felmer--Loss--Paturel~\cite{DFLP}. Hong--Kwon--Yoon
\cite{HKY} give fractional operator kinetic estimates and conservation
of pairwise component inner products. We do not claim either idea as
new. The Gaussian coefficient $\mathcal B_r$ above is an explicit
refinement retaining the scalar endpoint and the full occupation
spectrum; its optimality at fixed higher rank is not asserted.
Indeed, for $k$ equal occupations the Gaussian ratio tends to
$1/\mathrm G(r+2)$ as $k\to\infty$, whereas operator-norm kinetic
inequalities can yield mass intervals growing with rank. Neither
bound uniformly replaces the other for every purpose.

Frank--Gontier--Lewin~\cite{FGL} prove strict finite-rank binding for
the ordinary second-order kinetic problem in their specified range.
For eigenvalue moment one and dimension at least three, rank two
strictly improves on the rank-one spectral constant. Consequently
one cannot simply replace total mass by largest occupation while
retaining the scalar optimal kinetic constant. This is an established
obstruction from their work; we neither solve a Lieb--Thirring
conjecture nor transfer that strict-binding theorem to quartic
dispersion without proof.

\paragraph{What remains mathematically open in this manuscript.}
The mean bounds permit bounded global motion at a sufficiently large
kinetic level. They also permit sparse high peaks when no kinetic
upper bound is assumed. Applying the same estimate at a later time
does not amplify it again, because $d$ and $E$ are conserved.
To exclude every bounded negative-energy trajectory one needs a new
long-time statement controlling scale evolution or the accumulated
motion of the small mass tail. Only after that step could one
separate finite-time collapse from infinite-time growth. The
modulation identities locate these terms explicitly; they do not
remove them. Independently, proving real-sector nondegeneracy in
the full ground-wave class would remove the remaining hypothesis
from the stronger center theorem. Determining the optimal
fixed-occupation constant is another substantive variational problem.

\appendix
\section{Scalar compactness and critical normalization}
\label{sec:r10-foundations}
The Fourier transform is unitary, $D=-i\nabla$, and
$P(u)=\operatorname{Im}\int\overline u\nabla u=\int\xi|\widehat u|^2$.
The real Hilbert pairing is the real part of the complex $L^2$ pairing.
Throughout the local-power arguments, $N\ge1$, $m>1$, $s=m/2$,
$r=m/N$, and $p=2+2r$. A principal symbol $a$ is real, even,
$m$-homogeneous, $C^{N+2}$ off zero, and satisfies
$c|\xi|^m\le a(\xi)\le C|\xi|^m$.
Set $A(u)=\int a|\widehat u|^2$, $M(u)=\|u\|_2^2$, and
$E(u)=A(u)/2-\|u\|_p^p/p$. The sharp critical mass $M_c$ is defined by
\begin{equation}\label{eq:r10-basic-GN}
 \|u\|_p^p\le\frac p2\left(\frac{M(u)}{M_c}\right)^r A(u).
\end{equation}
The traveling-wave convention $u(t,x)=e^{i\omega t}U(x-vt)$ gives
$[a(D)+\omega+iv\cdot\nabla]U-|U|^{2r}U=0$.

\begin{lemma}[Normalized optimizer compactness]
The inequality \eqref{eq:r10-basic-GN} is attained. If
$M(f_n)\to M_c$, $A(f_n)\to1$, and $\|f_n\|_p^p\to p/2$, then
some translated subsequence converges strongly in $H^s$ to an optimizer.
The normalized optimizer set is compact modulo translations and phases.
\end{lemma}
\begin{proof}
Sharp Sobolev interpolation gives a finite positive best constant, and
amplitude and spatial scaling normalize a maximizing sequence as stated.
Bounded energy-space sequences that vanish locally in mass converge to
zero in $L^p$: use a bounded-overlap cover by unit balls, local Sobolev
interpolation, and $2<p<2N/(N-m)$ when $N>m$ (all finite $p$ otherwise).
Thus translations give a nonzero weak limit $Q$ and almost-everywhere
convergence. Write $M_1=M(Q)>0$, $A_1=A(Q)>0$, $M_2=M_c-M_1$,
and $A_2=1-A_1$. Quadratic splitting and Brezis--Lieb give
\[
 M_c^r\le M_1^r A_1+M_2^r A_2,\qquad A_1+A_2=1.
\]
If both masses are positive, the right side is strictly below $M_c^r$.
Consequently $M_1=M_c$ and then $A_1=1$. Weak convergence plus
convergence of both quadratic norms gives strong $H^s$ convergence.
The same proof applies to a sequence of optimizers. Constant phase is
proved by the density-norm lemma below; this compactness argument itself
does not need a positivity-preserving kinetic operator.
\end{proof}


\section{Sharp density geometry and quantitative polarization}
\label{sec:r12-polarization-foundation}
This appendix reproduces the analytic inputs used by the new evolution
arguments. They are proved here, including the infinite-dimensional
compactness step. No radiality, positivity of the kinetic semigroup,
or real-sector nondegeneracy is used.

Let $\mathcal H$ be any nonzero separable complex Hilbert space.
Finite-dimensional $\mathbb C^d$ and $\ell^2$ are both allowed.
Fix $N\ge1$, $m>1$, $s=m/2$, and a real-even, positive elliptic,
$m$-homogeneous symbol $a$, of class $C^{N+2}$ away from zero.
For $U\in H^s(\mathbb R^N;\mathcal H)$ put
\[
 \rho_U=\|U\|_{\mathcal H}^2,\quad
 M(U)=\int\rho_U,\quad
 A(U)=\int a(\xi)\|\widehat U(\xi)\|_{\mathcal H}^2\,d\xi.
\]
We use either of the following interactions:
\begin{align*}
 &\kappa=1+m/N,
 &&\mathcal F(U)=\int\rho_U^\kappa,
 &&V_\rho=\rho^{\kappa-1};\\
 &\kappa=2,
 &&\mathcal F(U)=\mathcal D(\rho_U,\rho_U),
 &&V_\rho=K*\rho.
\end{align*}
In the second case $1<m<N$, the real even kernel $K$ is homogeneous
of degree $-m$, its interaction form $\mathcal D$ is positive definite,
and
\[
 |K(x)|+|x||\nabla K(x)|\le C|x|^{-m}.
\]
As before, positivity of the Fourier transform of $K$ almost everywhere
suffices for positive definiteness. Pointwise positivity is not needed.
The scalar sharp inequality defines $M_c$ by
\begin{equation}\label{eq:r11-vectorGN}
 \mathcal F(u)\le\kappa
       \left(\frac{M(u)}{M_c}\right)^{\kappa-1}A(u).
\end{equation}

\begin{theorem}[Sharp interpolation and dimension-independent compactness]
\label{thm:r11-vectorcompact}
Inequality \eqref{eq:r11-vectorGN} has the same sharp constant for
$\mathcal H$-valued functions. Every optimizer has the form
\[
 U(x)=\zeta Q(x),\qquad \|\zeta\|_{\mathcal H}=1,
\]
where $Q$ is a real scalar optimizer, up to a constant phase.
If $M(U_n)\to M_c$, $A(U_n)\to1$, and
$\mathcal F(U_n)\to\kappa$, then the sequence is strongly compact
in $H^s$ after translations and internal unitaries. This remains
true when the internal dimensions vary: embed all spaces in
$\ell^2$ before applying the unitaries.
\end{theorem}

\subsection{Sharp interpolation and complete polarization}
The local interaction is the $\kappa$-th power of the $L^\kappa$
density norm. The Hartree interaction is the square of its positive
definite Hilbert density norm. Thus both have the form
$\mathcal F(U)=\|\rho_U\|_{\mathcal B}^\kappa$ for a strictly convex
real density space $\mathcal B$.
In an internal orthonormal basis $U=(u_j)$, the triangle inequality,
scalar sharp interpolation, and sequence H\"older give
\begin{align*}
 \mathcal F(U)^{1/\kappa}
 &\le\sum_j\mathcal F(u_j)^{1/\kappa}\\
 &\le\kappa^{1/\kappa}M_c^{-(\kappa-1)/\kappa}
   \sum_j A(u_j)^{1/\kappa}M(u_j)^{(\kappa-1)/\kappa}\\
 &\le\kappa^{1/\kappa}M_c^{-(\kappa-1)/\kappa}
              A(U)^{1/\kappa}M(U)^{(\kappa-1)/\kappa}.
\end{align*}
Finite partial sums converge in the density norm because the same
bound applied to their tails tends to zero. A scalar maximizing sequence
in one component proves sharpness.

Split the components further into real scalar functions $f_j$.
Equality forces every active $f_j$ to be a scalar optimizer, all
$f_j^2$ to be proportional, and all kinetic-to-mass ratios to agree.
The same conclusions hold after any real orthogonal rotation of a
pair $f_i,f_j$. If $f_j/f_i$ took both signs of a nonzero constant,
the squared sum and difference could not be proportional; when the
constant has magnitude one their nonzero supports would be disjoint.
Thus each $f_j$ is one real constant multiple of a fixed scalar
function. This proves complete polarization, including for countably
many components.

\subsection{Compactness modulo internal unitaries}
Ordinary local compactness does not hold directly in an infinite
internal dimension: a fixed spatial shape can move through mutually
orthogonal internal vectors. The following argument removes precisely
that loss of compactness.

\begin{lemma}[Internal-unitary local compactness]
\label{lem:r11-internalcompact}
Let $X=H^s(\mathbb R^N;\mathbb C)$ with $s>0$. From every bounded
sequence in $X\widehat\otimes\ell^2$, one can extract a subsequence
and internal unitaries $W_n$ such that $W_nU_n$ converges weakly in
$X\widehat\otimes\ell^2$ and strongly in
$L^2_{\mathrm{loc}}(\mathbb R^N;\ell^2)$.
\end{lemma}
\begin{proof}
Fix an orthonormal basis $(\phi_k)$ of the scalar Hilbert space $X$.
Write its internal coefficient vectors as
$c_{nk}=\langle\phi_k,U_n\rangle_X\in\ell^2$.
Gram--Schmidt the first $n$ vectors and extend the resulting finite
orthonormal list to an internal orthonormal basis. A unitary $W_n$
can therefore be chosen so that
\[
 W_nc_{nk}\in\operatorname{span}\{e_1,\ldots,e_k\},\qquad k\le n.
\]
For each fixed $k$, these vectors have a convergent subsequence in a
fixed finite-dimensional space. Diagonal extraction gives strong
convergence of every coefficient vector. Weak compactness gives a
weak limit of $W_nU_n$, with exactly those coefficients.

For $\chi\in C_c^\infty$, the scalar multiplication map
$K_\chi:X\to L^2$, $f\mapsto\chi f$, is compact. If $P_L$ is the
scalar projection onto the first $L$ basis vectors, then
\[
 \|K_\chi(I-P_L)\|_{X\to L^2}\longrightarrow0.
\]
Tensoring with the identity on $\ell^2$ preserves this norm.
The transformed $P_LU_n$ converge strongly because there are only
$L$ coefficient vectors; after multiplication by $\chi$, the tails
are uniformly small. This proves local strong convergence.
\end{proof}

\begin{proof}[Compactness assertion of Theorem~\ref{thm:r11-vectorcompact}]
The localized Sobolev inequality has constants independent of internal
dimension: use the Hilbert-valued Bessel-potential representation and
its norm triangle inequality, then spatial localization and the sum of
component quadratic norms. Thus local mass vanishing for an
$H^s$-bounded sequence implies vanishing in $L^p(\mathcal H)$ for
$2<p<2N/(N-2s)$, or every finite $p>2$ if $2s\ge N$.
In the local case take $p=2\kappa$; in the Hartree case use
$p=4N/(2N-m)$ and Hardy--Littlewood--Sobolev. Hence
$\mathcal F(U_n)\to\kappa>0$ excludes local mass vanishing.

Translate a nonzero amount of local mass into a fixed ball and use
Lemma~\ref{lem:r11-internalcompact}. We obtain a nonzero weak limit
$U$, strong local $L^2$ convergence, and by interpolation strong local
convergence in the interaction exponent. Put $R_n=U_n-U$.
The quadratic quantities split. For local powers the Brezis--Lieb
argument for $\|U\|_{\mathcal H}^{2\kappa}$ gives
\[
 \mathcal F(U_n)=\mathcal F(U)+\mathcal F(R_n)+o(1).
\]
For Hartree, the term $\operatorname{Re}\langle U,R_n\rangle$
tends to zero in $L^{p/2}$ by local convergence and the small
$L^p$ tail of $U$. Also
$\mathcal D(\rho_U,\rho_{R_n})\to0$: approximate $U$ in
$L^p(\mathcal H)$ by a compactly supported smooth finite-dimensional
field. Its potential is bounded and decays at infinity, while $R_n$
tends locally to zero in $L^2$. HLS controls the approximation error.
This proves the same splitting for Hartree.

Let $M_1>0,M_2$ and $A_1,A_2$ be the limiting contributions.
Then $M_1+M_2=M_c$, $A_1+A_2=1$, and sharp interpolation yields
\[
 M_c^{\kappa-1}\le
 M_1^{\kappa-1}A_1+M_2^{\kappa-1}A_2.
\]
Both masses cannot be positive, since each would be strictly below
$M_c$. Thus $M_2=0$, and the inequality then forces $A_2=0$.
The convergence is strong in $H^s$. The limit is an optimizer,
already proved to be completely polarized.
\end{proof}

\subsection{Phase coercivity and the internal symmetry kernel}
For completeness, the scalar phase fact used below follows directly
from the density norm. Normalize a real optimizer by $M(Q)=M_c$,
$A(Q)=1$, and put $\omega_*=(\kappa-1)/M_c$. The form
\[
 L_Q=a(D)+\omega_*-V_{Q^2}
\]
is nonnegative: differentiate the sharp inequality at
$Q+i\sqrt t\,h$. If $h$ belongs to its kernel, write
$x=A(h)$ and $y=M(h)/M_c$. The supporting-functional inequality
for the density norm, sharp interpolation for $h$, and weighted
arithmetic--geometric mean give
\[
 x+(\kappa-1)y
 \le\kappa x^{1/\kappa}y^{(\kappa-1)/\kappa}
 \le x+(\kappa-1)y.
\]
Equality implies $x=y$ and $h^2=yQ^2$ by strict convexity.
The kernel of a nonnegative form is linear; apply the same conclusion
to $h+Q$ to get $h=cQ$. The potential is compact as a form
$H^s\to H^{-s}$, by local Rellich and H\"older in the local case,
or the compact-potential approximation and HLS in the Hartree case.
Consequently
\begin{equation}\label{eq:r11-vectorcoercivity}
 \langle h,L_Qh\rangle\ge c_Q\|h\|_{H^s}^2,
 \qquad \langle h,Q\rangle_{L^2}=0.
\end{equation}
This holds for complex $h$ by real--imaginary splitting, and it can
be summed over arbitrarily many components with the same constant.

At $Qe_1$, write a perturbation as $(h+ik)e_1+Z$ with $Z\perp e_1$.
The Hessian of $E+\omega_*M/2$ is
\[
 \langle h,L_+h\rangle+\langle k,L_Qk\rangle
 +\sum_{j\ge2}\bigl(
  \langle\operatorname{Re}Z_j,L_Q\operatorname{Re}Z_j\rangle
 +\langle\operatorname{Im}Z_j,L_Q\operatorname{Im}Z_j\rangle\bigr).
\]
Its kernel in the phase and internal transverse sector is exactly
\[
 \{Q\eta:\ \eta\in\mathcal H,
                 \operatorname{Re}\langle\eta,e_1\rangle=0\}.
\]
For $\mathcal H=\mathbb C^d$, this is the $(2d-1)$-dimensional real
tangent space to the internal unitary orbit. Any additional full
Hessian kernel lies in the scalar real sector $L_+$; no assumption
on that sector is made.

\subsection{Quantitative polarization without real-sector nondegeneracy}
The preceding compactness and phase coercivity also give a quantitative
statement for arbitrary fields, not just stationary or traveling ones.
Define the scale-invariant interpolation deficit
\[
 \mathfrak d(U)=1-
 \frac{\mathcal F(U)}{\kappa(M(U)/M_c)^{\kappa-1}A(U)}
\]
and the squared distance to a completely polarized field by
\[
 \mathfrak p(U)=\inf_{\substack{\|\zeta\|_{\mathcal H}=1\\q\in H^s}}
 \left\{\frac{M(U-\zeta q)}{M(U)}
             +\frac{A(U-\zeta q)}{A(U)}\right\}.
\]
Both are defined for every nonzero $U\in H^s(\mathcal H)$.

\begin{theorem}[Dimension-independent quantitative polarization]
\label{thm:r11-quantitativepolarization}
There is $C<\infty$, depending only on the fixed scalar dispersion
and interaction, such that
\begin{equation}\label{eq:r11-polarizationdeficit}
 \mathfrak p(U)\le C\mathfrak d(U)
\end{equation}
for every nonzero Hilbert-valued field $U$ and every internal dimension.
No nondegeneracy or uniqueness of the scalar real-sector ground state
is required.
\end{theorem}
\begin{proof}
Amplitude and spatial scaling preserve both quantities, so impose
$M(U)=M_c$, $A(U)=1$. Write
$\|U\|_X^2=M(U)/M_c+A(U)$.
Since $\mathfrak p(U)\le2$, failure of a uniform constant would
produce $U_n$ with $\mathfrak d(U_n)\to0$ and
$\mathfrak d(U_n)/\mathfrak p(U_n)\to0$.
After the symmetries of Theorem~\ref{thm:r11-vectorcompact},
$U_n\to Qe_1$ strongly in $H^s$.
Align $\int QU_n$ with the positive real first internal direction
and write $U_n=(u_n,Z_n)$. Each component of $Z_n$ is complex
$L^2$-orthogonal to $Q$, and $u_n\to Q$, $Z_n\to0$ in $H^s$.

The interaction expansion in these transverse directions is
\begin{equation}\label{eq:r11-transverseexpansion}
 \mathcal F(u_n,Z_n)=\mathcal F(u_n)
 +\kappa\langle V_{|u_n|^2},\|Z_n\|^2\rangle
 +o(\|Z_n\|_X^2).
\end{equation}
For local powers, set $a=|u_n|^2$ and $b=\|Z_n\|^2$ pointwise.
When $1<\kappa\le2$, the nonnegative Taylor remainder of
$(a+b)^\kappa$ after its linear term in $b$ is at most $C_\kappa b^\kappa$.
When $\kappa>2$, it is at most
$C_\kappa(a^{\kappa-2}b^2+b^\kappa)$.
H\"older and the dimension-independent $L^{2\kappa}$ Sobolev bound
make these respectively $O(\|Z_n\|_X^{2\kappa})$ and
$O(\|Z_n\|_X^4+\|Z_n\|_X^{2\kappa})$, since $u_n$ stays bounded.
Both are $o(\|Z_n\|_X^2)$.
For Hartree the remainder is exactly
$\mathcal D(\|Z_n\|^2,\|Z_n\|^2)=O(\|Z_n\|_X^4)$.

Put $m_n=M(Z_n)$ and $a_n=A(Z_n)$. Scalar sharp interpolation gives
\[
 \mathcal F(u_n)\le
 \kappa(1-m_n/M_c)^{\kappa-1}(1-a_n).
\]
Combining its expansion with \eqref{eq:r11-transverseexpansion}
and potential form convergence yields
\begin{align*}
 \kappa-\mathcal F(U_n)
 &\ge\kappa\left[
 a_n+\frac{\kappa-1}{M_c}m_n
       -\langle V_{Q^2},\|Z_n\|^2\rangle\right]
       -o(\|Z_n\|_X^2)\\
 &=\kappa\sum_j\langle Z_{nj},L_QZ_{nj}\rangle
       -o(\|Z_n\|_X^2)\\
 &\ge c\|Z_n\|_X^2.
\end{align*}
The constant is independent of component number. Taking the candidate
$e_1u_n$ in the distance gives
$\mathfrak p(U_n)\le\|Z_n\|_X^2$.
As $\mathfrak d(U_n)=(\kappa-\mathcal F(U_n))/\kappa$,
this contradicts the assumed ratio and proves the theorem.
\end{proof}

The theorem controls internal polarization rather than the full
scalar shape. It therefore remains valid even when scalar optimizer
uniqueness or real-sector quantitative stability is unavailable.

\subsection{The real version and scalar constant phase}
\label{subsec:r13-real-phase}
All preceding density inequalities, compactness, and quantitative
polarization statements hold over a separable real internal Hilbert
space, with orthogonal maps replacing unitaries and real scalar $q$
in the definition of $\mathfrak p$. Here are the points that matter
for the scalar evolution theorem. The real scalar sharp constant
equals the complex one: apply the density triangle inequality and
sequence H\"older to $(\operatorname{Re}u,\operatorname{Im}u)$,
starting with the real scalar constant, while the reverse comparison
is immediate. Equality forces a constant real internal direction
by the pair-rotation argument above. The internal compactness proof
uses real Gram--Schmidt without any other change. Finally rotate
$\int QU_n$ into the first real direction. Every transverse real
component is orthogonal to $Q$, so exactly the displayed $L_Q$
coercivity and transverse expansion apply. This proves the real
quantitative assertion, not just its complex analogue.

In particular, for a scalar complex $u$ and either allowed interaction,
there exists a phase $\theta$ such that $u=e^{i\theta}(X+iY)$,
with $X,Y$ real, and
\[
 \frac{M(Y)}{M(u)}+\frac{A(Y)}{A(u)}\le C\mathfrak d(u).
\]
This holds for every local exponent $r>0$, and for the strictly
positive-definite Hartree density form. For negative energy and
$d=(M/M_c)^{\kappa-1}-1>0$, one has
$0\le\mathfrak d(u)<d/(1+d)$, providing exactly the input used
in Theorem~\ref{thm:r13-general-average}. No phase is differentiated
in time. In the local range $r\ge1$,
Theorem~\ref{thm:r13-explicitpolarization} gives an explicit constant
in place of this compactness argument.

\section{Zero momentum and instantaneous quartic transport}
The distinction between momentum and quartic mass transport is exact.
Choose a real sharp wave $Q$ and a real smooth test $h$ with
\[
 \langle h,Q\rangle=0,\qquad
 \langle h,\partial_kQ\rangle=0\ (1\le k\le N),\qquad
 \langle h,\partial_j\Delta Q\rangle\ne0
\]
for some $j$. Such a test exists. Otherwise
$\partial_j\Delta Q\in\operatorname{span}\{Q,\partial_1Q,\ldots,
\partial_NQ\}$ as distributions. Pairing with $Q$ removes its $Q$
component, and Fourier transformation would force $\widehat Q$ onto
the zero set of a nonzero polynomial, contradicting $Q\ne0$ in $L^2$.
For fixed sufficiently large $a>0$, set
\[
 u_\varepsilon=(1+a\varepsilon^2)Q+i\varepsilon h.
\]
Then $P(u_\varepsilon)=0$ exactly, and
\begin{align*}
 d_\varepsilon
 &=\frac4N\left(2a+\frac{M(h)}{M_c}\right)\varepsilon^2
       +O(\varepsilon^4),\\
 E(u_\varepsilon)
 &=\varepsilon^2\left[
 \frac{A(h)-\int|Q|^{8/N}h^2}{2}-a\omega M_c\right]
       +o(\varepsilon^2)<0,
\end{align*}
where $\Delta^2Q+\omega Q-|Q|^{8/N}Q=0$.
Nevertheless its quartic mass current is
\[
 4\int\xi_j|\xi|^2|\widehat u_\varepsilon(\xi)|^2\,d\xi
 =8\varepsilon(1+a\varepsilon^2)
             \langle h,\partial_j\Delta Q\rangle.
\]
Thus this current has order $\sqrt{d_\varepsilon}$ even at exactly zero
momentum and negative energy. With a finite first moment it is the
instantaneous derivative of that moment; without one it remains the
finite Fourier current underlying localized mass transport.

The improved center in Section~\ref{sec:r13-center} is a modulation
coordinate chosen by dual orthogonality, not the ordinary first-moment
center. Its quadratic drift is therefore consistent with the example.
The example rules out the proposed instantaneous current improvement;
it does not rule out further cancellation along the nonlinear evolution.


\clearpage
\begin{thebibliography}{9}
\raggedright
\bibitem{BL} T. Boulenger and E. Lenzmann,
\emph{Blowup for biharmonic NLS}, Ann. Sci. ENS 50 (2017).
\href{https://arxiv.org/abs/1503.01741}{arXiv:1503.01741}.
\bibitem{BCGJ} D. Bonheure, J.-B. Cast\'eras, T. Gou and L. Jeanjean,
\emph{Strong instability of ground states to a fourth order Schr\"odinger equation},
IMRN 2019, 5299--5315.
\href{https://arxiv.org/abs/1703.07977}{arXiv:1703.07977}.
\bibitem{DM} L. Dinh, with Muse Spark,
\emph{Finite-Time Blow-Up of Radial Negative-Energy Solutions for the
Mass-Critical Biharmonic Nonlinear Schr\"odinger Equation}, manuscript,
2 October 2026.
\href{https://ai.meta.com/research/publications/finite-time-blow-up-of-radial-negative-energy-solutions-for-the-mass-critical-biharmonic-nonlinear-schrodinger-equation/}{Official manuscript and PDF}.
\bibitem{PS} B. Pausader and S. Shao,
\emph{The mass-critical fourth-order Schr\"odinger equation in high dimensions},
J. Hyperbolic Differ. Equ. 7 (2010), 651--705.
\href{https://arxiv.org/abs/0906.1547}{arXiv:0906.1547}.
\bibitem{IKN} T. Inui, K. Komada and K. Nishimura,
\emph{Scattering of non-radial solutions for the fourth-order
Schr\"odinger equation with $\dot H^{1/2}$-critical and supercritical nonlinearity},
2026 preprint.
\href{https://arxiv.org/abs/2610.05802}{arXiv:2610.05802v1}.
\bibitem{DFLP} J. Dolbeault, P. Felmer, M. Loss and E. Paturel,
\emph{Lieb--Thirring type inequalities and Gagliardo--Nirenberg
inequalities for systems}, J. Funct. Anal. 238 (2006), 193--220.
\href{https://arxiv.org/abs/math-ph/0506052}{arXiv:math-ph/0506052}.
\bibitem{HKY} Y. Hong, S. Kwon and H. Yoon,
\emph{Global existence versus finite time blowup dichotomy for the system
of nonlinear Schr\"odinger equations}.
\href{https://arxiv.org/abs/1712.06811}{arXiv:1712.06811}.
\bibitem{FGL} R. L. Frank, D. Gontier and M. Lewin,
\emph{The Nonlinear Schr\"odinger Equation for Orthonormal Functions II:
Application to Lieb--Thirring Inequalities},
Commun. Math. Phys. 384 (2021), 1783--1828.
\href{https://arxiv.org/abs/2002.04964}{arXiv:2002.04964}.
\end{thebibliography}
\end{document}
